% 第1章 はじめに。oshifit_spec.md 第1・2節にもとづく下書き。【】の赤字は未確認・未記入。
\chapter{はじめに}

\section{研究の背景}
アイドルやアニメのキャラクター、VTuber などを応援する「推し活」は、若い世代を中心に広く行われている\todo{要出典：推し活の普及に関する調査}。
推し活は、生きがいや日々のストレスの解消、人とのつながりをもたらす一方で、睡眠時間を削る、予算を超えてお金を使う、イベントや情報を追うことが義務のように感じられる、SNS で他のファンと比べて落ち込む、といった負担（以下、推し疲れ）につながることもある\todo{要出典：推し活の効果と負担}。

筆者らが大学生を中心とする215名に行った予備調査でも、推し活の頻度は、生きがいや自信といった良い面と、睡眠や金銭の不安といった負担の両方と関連していた。
推し活は「多いほど良い」「少ないほど良い」と一方向には言えず、その人にとってちょうどよい形を探すことが大切だと考えられる。

しかし、推し活の負担に対する従来の助言は「使うお金を減らす」「時間を減らす」といった、推し活そのものを控える方向に偏りがちである\todo{要出典}。
推し活は本人にとって大切な楽しみであり、控えること自体が満足感を下げるおそれがある。
推し活の楽しさを保ったまま、負担を生む「やり方」だけを見直す支援が求められる。

\section{研究の目的}
本研究では、推し活の状態を5つの要素（推し活の充実、睡眠と生活、お金、義務感と人間関係、SNSと比較）で見える化し、本人の性格と活動の比重に合わせた小さな提案を行うシステム OshiFit を設計・実装する。
OshiFit は、提案を取り入れた場合（プランB）と今のまま続けた場合（プランA）の4週間後の見込みを、本人に似せたアバターで並べて示し、どちらを選ぶかを本人に委ねる。
また、推し活がすでに順調な人には、提案をしない。

そのうえで、4週間の使用を通じて、次の2点を検討する。
\begin{enumerate}
  \item 提示した見込みが、実際の変化とどの程度一致したか（提案の精度）。
  \item 提案を選んで実行した人で、幸福度とストレスがどのように変化したか（心理的な効果）。
\end{enumerate}

\section{本研究の意義}
本研究の意義は、推し活を支援するシステムを作るだけでなく、その人に合った推し活の提案が、幸福度の向上やストレスの緩和といったウェルビーイングと結びつくかを、実際の使用データで確かめる点にある。
具体的には、次の点に新しさがある。
\begin{itemize}
  \item 推し活を「減らす」のではなく、5つの要素と性格・活動の比重にもとづいて「やり方を組み替える」提案を行う。
  \item 提案ごとに、根拠となる研究と、その確かさを3段階で示す。
  \item 本人に似せたアバターで「2つの未来」を並べ、選ぶのは本人に委ねる。
  \item 推し活が順調な人には「今のままで順調です」と伝え、提案をしない。
\end{itemize}
\todo{独自性の最終的な言い方は研究者が確定する（oshifit\_spec.md 第2節）}

なお、本研究では、提案を受けるかどうかを本人が選ぶため、群の違いには本人の特徴の違いも含まれる。
したがって、本研究の結果は提案の効果を断定するものではなく、提案を選んで実行することと、幸福度・ストレスの変化との関連を示すものである。

\section{本論文の構成}
第2章では関連研究を、第3章では OshiFit の設計を述べる。
第4章で実験の方法を、第5章で結果を示し、第6章で考察と今後の課題を述べる。
